% LaTeX Template For MATH 490 @ VCU
\documentclass{article}

\usepackage{hyperref}
\usepackage{amsmath}
\usepackage{amsthm}
\usepackage{amssymb}
\usepackage{enumerate}

\DeclareMathOperator{\arcsec}{arcsec}
\DeclareMathOperator{\arccot}{arccot}
\DeclareMathOperator{\arccsc}{arccsc}

\linespread{2}
\setlength{\textwidth}{6.5in}
\setlength{\textheight}{8.5in}
\setlength{\oddsidemargin}{0in}
\setlength{\evensidemargin}{0in}
\setlength{\topmargin}{0in}
\setlength{\headheight}{0in}
\setlength{\headsep}{0in}
\setlength{\footskip}{.7in}

\makeatletter
\newcommand{\vast}{\bBigg@{4}}
\newcommand{\Vast}{\bBigg@{5}}
\makeatother

\begin{document}

\title{National Yang Ming Chiao Tung University\\Calculus A (I) Report 2}
\author{Chang, Yung-Hsuan\\111652004\\Department of Applied Mathematics}
\maketitle
\newpage

\section*{Question 1}
\begin{enumerate}[(1)]
    \item Compute the integral $\displaystyle\int_0^\infty{\dfrac{\arctan{ax}-\arctan{bx}}{x}dx}$, where $a<b$.\\
    \textbf{Solution:}\\
    This is an improper integral. Let $I=\displaystyle\int_0^\infty{\dfrac{\arctan{ax}-\arctan{bx}}{x}dx}$, then we have
    \begin{align*}
        I=&\int_0^1{\dfrac{\arctan{ax}-\arctan{bx}}{x}dx}+\int_1^\infty{\dfrac{\arctan{ax}-\arctan{bx}}{x}dx}\\
        =&\lim_{\varepsilon\to0^+}{\int_\varepsilon^1{\dfrac{\arctan{ax}-\arctan{bx}}{x}dx}}+\lim_{R\to\infty}{\int_1^R{\dfrac{\arctan{ax}-\arctan{bx}}{x}dx}}.
    \end{align*}
    We discuss for five cases below:
    \begin{enumerate}
        \item $0<a<b$\label{1.(a)}\\
        In this case, we have
        \begin{align*}
            I&=\lim_{\varepsilon\to0^+}{\int_\varepsilon^1{\dfrac{\arctan{ax}-\arctan{bx}}{x}dx}}+\lim_{R\to\infty}{\int_1^R{\dfrac{\arctan{ax}-\arctan{bx}}{x}dx}}\\
            &=\lim_{\varepsilon\to0^+}{\int_\varepsilon^1{\dfrac{\arctan{ax}}{x}dx}}-\lim_{\varepsilon\to0^+}{\int_\varepsilon^1{\dfrac{\arctan{bx}}{x}dx}}\\
            &\quad\:+\lim_{R\to\infty}{\int_1^R{\dfrac{\arctan{ax}}{x}dx}}-\lim_{R\to\infty}{\int_1^R{\dfrac{\arctan{bx}}{x}dx}}.
        \end{align*}
        Let $u=ax$ in the first and third term; let $u=bx$ in the second and fourth term, then we have
        \begin{align*}
            I&=\lim_{\varepsilon\to0^+}{\int_{a\varepsilon}^a{\dfrac{\arctan{u}}{\dfrac{u}{a}}\dfrac{1}{a}du}}-\lim_{\varepsilon\to0^+}{\int_{b\varepsilon}^b{\dfrac{\arctan{u}}{\dfrac{u}{b}}\dfrac{1}{b}du}}\\
            &\quad\:+\lim_{R\to\infty}{\int_a^{aR}{\dfrac{\arctan{u}}{\dfrac{u}{a}}\dfrac{1}{a}du}}-\lim_{R\to\infty}{\int_b^{bR}{\dfrac{\arctan{u}}{\dfrac{u}{b}}\dfrac{1}{b}du}}\\
            &=\lim_{\varepsilon\to0^+}{\int_{a\varepsilon}^a{\dfrac{\arctan{u}}{u}du}}-\lim_{\varepsilon\to0^+}{\int_{b\varepsilon}^b{\dfrac{\arctan{u}}{u}du}}\\
            &\quad\:+\lim_{R\to\infty}{\int_a^{aR}{\dfrac{\arctan{u}}{u}du}}-\lim_{R\to\infty}{\int_b^{bR}{\dfrac{\arctan{u}}{u}du}}\\
            &=\lim_{\varepsilon\to0^+}{\int_{a\varepsilon}^{b\varepsilon}{\dfrac{\arctan{u}}{u}du}}-\lim_{\varepsilon\to0^+}{\int_a^b{\dfrac{\arctan{u}}{u}du}}\\
            &\quad\:+\lim_{R\to\infty}{\int_{aR}^{bR}{\dfrac{\arctan{u}}{u}du}}-\lim_{R\to\infty}{\int_b^a{\dfrac{\arctan{u}}{u}du}}.\\
        \end{align*}
        Since the second and fourth term cancel out, we have $$I=\lim_{\varepsilon\to0^+}{\int_{a\varepsilon}^{b\varepsilon}{\dfrac{\arctan{u}}{u}du}}+\lim_{R\to\infty}{\int_{aR}^{bR}{\dfrac{\arctan{u}}{u}du}}.$$
        By M.V.T., there exists $cr\in(ar,br)$ for $r$ in $\mathbb{R}^+$ such that $$\int_{ar}^{br}{\dfrac{\arctan{u}}{u}du}=\arctan{cr}\int_{ar}^{br}{\dfrac{1}{u}du}.$$
        Thus, we have
        \begin{align*}
            I&=\lim_{\varepsilon\to0^+}{\arctan{c_1\varepsilon}\int_{a\varepsilon}^{b\varepsilon}{\dfrac{1}{u}du}}+\lim_{R\to\infty}{\arctan{c_2R}\int_{aR}^{bR}{\dfrac{1}{u}du}}\\
            &=\lim_{\varepsilon\to0^+}{\arctan{c_1\varepsilon}\ln{\dfrac{b\varepsilon}{a\varepsilon}}}+\lim_{R\to\infty}{\arctan{c_2R}\ln{\dfrac{aR}{bR}}}\\
            &=\ln{\dfrac{b}{a}}\lim_{\varepsilon\to0^+}{\arctan{c_1\varepsilon}}+\ln{\dfrac{a}{b}}\lim_{R\to\infty}{\arctan{c_2R}}.
        \end{align*}
        Since $c_1>0$ and $c_2>0$,
        \begin{align*}
            I&=\ln{\dfrac{b}{a}}\cdot0+\ln{\dfrac{a}{b}}\cdot\dfrac{\pi}{2}\\
            &=\dfrac{\pi}{2}\ln{\dfrac{a}{b}}.
        \end{align*}
        Therefore, $I=\dfrac{\pi}{2}\ln{\dfrac{a}{b}}$ whenever $0<a<b$.
        \item $0=a<b$\label{1.(b)}\\
        In this case, we have $$\displaystyle I=-\Big(\lim_{\varepsilon\to0^+}{\int_\varepsilon^1{\dfrac{\arctan{bx}}{x}dx}}+\lim_{R\to\infty}{\int_1^R{\dfrac{\arctan{bx}}{x}dx}}\Big).$$
        We observe that since $\dfrac{\arctan{x}}{x}\geq\dfrac{\arctan{1}}{x}$ for $x\geq1$, $\dfrac{\arctan{bx}}{x}\geq\dfrac{\arctan{b}}{x}$ for $b>0$ and $x\geq1$.\\
        Thus, we have $$\displaystyle\lim_{R\to\infty}{\int_1^R{\dfrac{\arctan{bx}}{x}dx}}\geq\lim_{R\to\infty}{\int_1^R{\dfrac{\arctan{b}}{x}dx}}.$$
        Since the limit $\displaystyle\lim_{R\to\infty}{\Big((\arctan{b})\int_1^R{\dfrac{1}{x}dx}}\Big)$ diverges, $\displaystyle\lim_{R\to\infty}{\int_1^R{\dfrac{\arctan{bx}}{x}dx}}$ diverges as well by the comparison test for improper integrals.\\
        Therefore, $I$ diverges whenever $0=a<b$.
        \item $a<0<b$\\
        Let $a=-\alpha$. Since arctangent is an odd function, we have $$\displaystyle I=-\Big(\lim_{\varepsilon\to0^+}{\int_\varepsilon^1{\dfrac{\arctan{\alpha x}+\arctan{bx}}{x}dx}}+\lim_{R\to\infty}{\int_1^R{\dfrac{\arctan{\alpha x}+\arctan{bx}}{x}dx}}\Big).$$
        We look at the second term: $$\lim_{R\to\infty}{\int_1^R{\dfrac{\arctan{\alpha x}+\arctan{bx}}{x}dx}}=\lim_{R\to\infty}{\int_1^R{\dfrac{\arctan{\alpha x}}{x}dx}}+\lim_{R\to\infty}{\int_1^R{\dfrac{\arctan{bx}}{x}dx}}.$$
        Since $\displaystyle\lim_{R\to\infty}{\int_1^R{\dfrac{\arctan{\alpha x}}{x}dx}}$ and $\displaystyle\lim_{R\to\infty}{\int_1^R{\dfrac{\arctan{bx}}{x}dx}}$ both diverge (we can prove these by using the same method with regard to the comparison test for improper integrals in \hyperref[1.(b)]{Question 1 (1) (b)}), $\displaystyle\lim_{R\to\infty}{\int_1^R{\dfrac{\arctan{\alpha x}+\arctan{bx}}{x}dx}}$ diverges; thus, $I$ diverges.\\
        Therefore, $I$ diverges whenever $a<0<b$.
        \item $a<b=0$\\
        Let $a=-\alpha$. Since arctangent is an odd function, we have $$\displaystyle I=-\Big(\lim_{\varepsilon\to0^+}{\int_\varepsilon^1{\dfrac{\arctan{\alpha x}}{x}dx}}+\lim_{R\to\infty}{\int_1^R{\dfrac{\arctan{\alpha x}}{x}dx}}\Big).$$
        Since $\displaystyle\lim_{R\to\infty}{\int_1^R{\dfrac{\arctan{\alpha x}}{x}dx}}$ diverges (we can prove this by using the same method with regard to the comparison test for improper integrals in \hyperref[1.(b)]{Question 1 (1) (b)}), $I$ diverges.\\
        Therefore, $I$ diverges whenever $a<b=0$.
        \item $a<b<0$\\
        Let $a=-\alpha$ and $b=-\beta$, then $\beta<\alpha$. Since arctangent is an odd function, we have
        \begin{align*}
            I&=-\Big(\lim_{\varepsilon\to0^+}{\int_\varepsilon^1{\dfrac{\arctan{\alpha x}-\arctan{\beta x}}{x}dx}}+\lim_{R\to\infty}{\int_1^R{\dfrac{\arctan{\alpha x}-\arctan{\beta x}}{x}dx}}\Big)\\
            \Leftrightarrow\:\:\:I&=\lim_{\varepsilon\to0^+}{\int_\varepsilon^1{\dfrac{\arctan{\beta x}-\arctan{\alpha x}}{x}dx}}+\lim_{R\to\infty}{\int_1^R{\dfrac{\arctan{\beta x}-\arctan{\alpha x}}{x}dx}}
        \end{align*}
        We find the answer just by algebraically replacing $\beta$ to the position of $a$ and replacing $\alpha$ to the position of $b$ in the result from \hyperref[1.(a)]{Question 1 (1) (a)}, having
        $$\lim_{\varepsilon\to0^+}{\int_\varepsilon^1{\dfrac{\arctan{\beta x}-\arctan{\alpha x}}{x}dx}}+\lim_{R\to\infty}{\int_1^R{\dfrac{\arctan{\beta x}-\arctan{\alpha x}}{x}dx}}=\dfrac{\pi}{2}\ln{\dfrac{\beta}{\alpha}}.$$
        Therefore, $I=\dfrac{\pi}{2}\ln{\dfrac{b}{a}}$ whenever $a<b<0$.
    \end{enumerate}
    \item Compute the limit $\displaystyle\lim_{n\to\infty}{\sum_{k=1}^{2n}{(-1)^k\Big(\dfrac{k}{2n}\Big)^{111}}}$.\\
    \textbf{Solution:}\\
    We first change the form of the summation, having 
    \begin{align*}
        \lim_{n\to\infty}{\sum_{k=1}^{2n}{(-1)^k\Big(\dfrac{k}{2n}\Big)^{111}}}&=\lim_{n\to\infty}{\Bigg(\quad(-1)\Big(\dfrac{1}{2n}\Big)^{111}+(-1)^3\Big(\dfrac{3}{2n}\Big)^{111}+\cdots+(-1)^{2n-1}\Big(\dfrac{2n-1}{2n}\Big)^{111}}\\
        &\quad\quad\quad\quad\:+(-1)^2\Big(\dfrac{2}{2n}\Big)^{111}+(-1)^4\Big(\dfrac{4}{2n}\Big)^{111}+\cdots+(-1)^{2n}\Big(\dfrac{2n}{2n}\Big)^{111}\Bigg)\\
        &=\lim_{n\to\infty}{\sum_{i=1}^{n}{(-1)\Big(\dfrac{2i-1}{2n}\Big)^{111}+\Big(\dfrac{2i}{2n}\Big)^{111}}}\\
        &=\lim_{n\to\infty}{\sum_{i=1}^{n}{\Big(\dfrac{2i}{2n}\Big)^{111}-\Big(\dfrac{2i-1}{2n}\Big)^{111}}}.
    \end{align*}
    Consider $f(x)=x^{111}$, by M.V.T., there exists $c_i$ in $\Big(\dfrac{2i-1}{2n}, \dfrac{2i}{2n}\Big)$ such that
    \begin{align*}
        \Big(\dfrac{d}{dx}f(x)\Big)\Big|_{x=c_i}&=\dfrac{f\Big(\dfrac{2i}{2n}\Big)-f\Big(\dfrac{2i-1}{2n}\Big)}{\dfrac{2i}{2n}-\dfrac{2i-1}{2n}}\\
        \Leftrightarrow\:\:\:\dfrac{1}{2n}\cdot\Big(\Big(\dfrac{d}{dx}f(x)\Big)\Big|_{x=c_i}\Big)&=f\Big(\dfrac{2i}{2n}\Big)-f\Big(\dfrac{2i-1}{2n}\Big).
    \end{align*}
    Thus, we have
    \begin{align*}
        \lim_{n\to\infty}{\sum_{i=1}^{n}{\Big(\dfrac{2i}{2n}\Big)^{111}-\Big(\dfrac{2i-2}{2n}\Big)^{111}}}&=\lim_{n\to\infty}{\sum_{i=1}^{n}{\dfrac{1}{2n}\cdot\Big(\Big(\dfrac{d}{dx}f(x)\Big)\Big|_{x=c_i}\Big)}}\\
        &=\dfrac{1}{2}\lim_{n\to\infty}{\sum_{i=1}^{n}{\Big(\Big(\dfrac{d}{dx}f(x)\Big)\Big|_{x=c_i}\Big)\dfrac{1}{n}}}\\
        &=\dfrac{1}{2}\int_0^1{\dfrac{d}{dx}f(t)dt}\\
        &=\dfrac{1}{2}\Big(f(x)\Big|^1_0\Big)\\
        &=\dfrac{1}{2}\big(1^{111}-0^{111}\big)\\
        &=\dfrac{1}{2}.
    \end{align*}
    Therefore, $\displaystyle\lim_{n\to\infty}{\sum_{k=1}^{2n}{(-1)^k\Big(\dfrac{k}{2n}\Big)^{111}}}=\dfrac{1}{2}$.
\end{enumerate}
\section*{Question 2}
We consider the function $\displaystyle f_n(x)=\dfrac{1}{1+x^2}-\sum_{k=1}^n{(-1)^{k-1}x^{2(k-1)}}$.
\begin{enumerate}[(1)]
    \item Show that $\displaystyle\lim_{n\to\infty}{\int_0^1{f_n(x)dx}}=0$.\\
    \textbf{Solution:}\\
    We directly substitute $\displaystyle f_n(x)=\dfrac{1}{1+x^2}-\sum_{k=1}^n{(-1)^{k-1}x^{2(k-1)}}$, having
    \begin{align*}
        \lim_{n\to\infty}{\int_0^1{f_n(x)dx}}&=\lim_{n\to\infty}{\int_0^1{\dfrac{1}{1+x^2}-\sum_{k=1}^n{(-1)^{k-1}x^{2(k-1)}}dx}}\\
        &=\lim_{n\to\infty}{\int_0^1{\dfrac{1}{1+x^2}-\sum_{k=1}^n{(-x^2)^{k-1}}dx}}\\
        &=\lim_{n\to\infty}{\int_0^1{\dfrac{1}{1+x^2}-1\cdot\dfrac{1-(-x^2)^n}{1-(-x^2)}}dx}\\
        &=\lim_{n\to\infty}{\int_0^1{\dfrac{(-x^2)^n}{1+x^2}dx}}.
    \end{align*}
    Let $\displaystyle I_n=\int_0^1{\dfrac{(-x^2)^n}{1+x^2}dx}$. Thus, we have
    \begin{align*}
        |I_n|&\leq\int_0^1{\dfrac{|(-x^2)^n|}{1+x^2}dx}\\
        &=\int_0^1{\dfrac{x^{2n}}{1+x^2}dx}\\
        &\leq\int_0^1{\dfrac{x^{2n}}{1}dx}\\
        &=\int_0^1{x^{2n}dx}\\
        &=\dfrac{1}{2n+1}x^{2n+1}\Big|_0^1\\
        &=\dfrac{1}{2n+1}\\
        \Rightarrow\:\:\:\lim_{n\to\infty}{|I_n|}&\leq\lim_{n\to\infty}\dfrac{1}{2n+1}\\
        &=0\\
        \Rightarrow\:\:\:\lim_{n\to\infty}{\;I_n\;}&=0.
    \end{align*}
    Therefore, $\displaystyle\lim_{n\to\infty}{\int_0^1{f_n(x)dx}}=0$.
    \item Using (1), show that $\displaystyle\sum_{n=1}^\infty{\dfrac{(-1)^{n-1}}{2n-1}}=\dfrac{\pi}{4}$.\\
    \textbf{Solution:}\\
    By $\displaystyle\lim_{n\to\infty}{\int_0^1{f_n(x)dx}}=0$, we have
    \begin{align*}
        0&=\lim_{n\to\infty}{\int_0^1{\dfrac{1}{1+x^2}-\sum_{k=1}^n{(-1)^{k-1}x^{2(k-1)}}dx}}\\
        &=\int_0^1{\dfrac{1}{1+x^2}dx}-\lim_{n\to\infty}{\int_0^1{\sum_{k=1}^n{(-1)^{k-1}x^{2(k-1)}}dx}}\\
        &=\arctan{x}\Big|_0^1-\lim_{n\to\infty}{\int_0^1{\sum_{k=1}^n{(-1)^{k-1}x^{2(k-1)}}dx}}\\
        &=(\arctan{1}-\arctan{0})-\lim_{n\to\infty}{\int_0^1{\sum_{k=1}^n{(-1)^{k-1}x^{2(k-1)}}dx}}\\
        &=\dfrac{\pi}{4}-\lim_{n\to\infty}{\int_0^1{\sum_{k=1}^n{(-1)^{k-1}x^{2(k-1)}}dx}}\\
        \Leftrightarrow\:\:\:\dfrac{\pi}{4}&=\lim_{n\to\infty}{\int_0^1{\sum_{k=1}^n{(-1)^{k-1}x^{2(k-1)}}dx}}.
    \end{align*}
    Since the integration is with respect to $x$, we have
    \begin{align*}
        \dfrac{\pi}{4}&=\lim_{n\to\infty}{\int_0^1{\sum_{k=1}^n{(-1)^{k-1}x^{2(k-1)}}dx}}\\
        &=\lim_{n\to\infty}{\sum_{k=1}^n{(-1)^{k-1}\int_0^1{x^{2(k-1)}dx}}}\\
        &=\lim_{n\to\infty}{\sum_{k=1}^n{(-1)^{k-1}\int_0^1{x^{2k-2}dx}}}\\
        &=\lim_{n\to\infty}{\sum_{k=1}^n{(-1)^{k-1}{\dfrac{1}{2k-1}x^{2k-1}\Big|^{x=1}_{x=0}}}}\\
        &=\lim_{n\to\infty}{\sum_{k=1}^n{\dfrac{(-1)^{k-1}}{2k-1}}}\\
        &=\sum_{k=1}^\infty{\dfrac{(-1)^{n-1}}{2n-1}}.
    \end{align*}
    Therefore, $\displaystyle\sum_{n=1}^\infty{\dfrac{(-1)^{n-1}}{2n-1}}=\dfrac{\pi}{4}$.
\end{enumerate}

\end{document}